\section{Particle filtering under known model assumptions}
A particle filter approximates sequential distributions in a specified state-space model. Let $x_t$ be the latent state at date $t$ and $y_t$ the observation then available. The transition density is $f(x_t\mid x_{t-1})$, the observation likelihood is $g(y_t\mid x_t)$, and a proposal $q$ generates the next particle. First consider a sequential importance-sampling step that extends particle $i$ without resampling it. Its updated weight satisfies
\begin{equation}
 w_t^{(i)}\propto w_{t-1}^{(i)}\frac{g(y_t\mid x_t^{(i)})f(x_t^{(i)}\mid x_{t-1}^{(i)})}{q(x_t^{(i)}\mid x_{t-1}^{(i)},y_t)}.
\end{equation}
The fraction is the incremental weight; multiplying it by the previous weight gives the new weight. It follows by dividing the joint posterior update by the proposal update, as in \citet[Section 3.3, equations (29)--(30)]{doucet2012tutorial}. The proposal must cover the relevant target support. A proposal that never visits a feasible high-risk region cannot recover it by weighting.

Resampling changes this calculation because selecting a parent already uses its weight. Let $W_{t-1}^{(j)}$ be the normalized weights before selection. Suppose child $i$ selects ancestor $A_t^{(i)}=j$ with probability $r_{t-1}^{(j)}$, then draws its next state from $q(\,\cdot\mid x_{t-1}^{(j)},y_t)$. Relative to this finite weighted approximation, the joint target mass for ancestor $j$ and new state $x$ is proportional to $W_{t-1}^{(j)} f(x\mid x_{t-1}^{(j)})g(y_t\mid x)$; the proposal mass is $r_{t-1}^{(j)}q(x\mid x_{t-1}^{(j)},y_t)$. Their ratio therefore gives
\begin{equation}
 \widetilde w_t^{(i)}\propto
 \frac{W_{t-1}^{(A_t^{(i)})}}{r_{t-1}^{(A_t^{(i)})}}
 \frac{g(y_t\mid x_t^{(i)})f(x_t^{(i)}\mid x_{t-1}^{(A_t^{(i)})})}
 {q(x_t^{(i)}\mid x_{t-1}^{(A_t^{(i)})},y_t)}.
\end{equation}
Both ancestor and state proposals must have positive mass wherever this target has positive mass. For ordinary multinomial resampling, $r_{t-1}^{(j)}=W_{t-1}^{(j)}$, so the first ratio cancels. Equivalently, the resampled children start with equal weights $1/N$ before propagation; the old parent's weight is not multiplied again. The always-resampled and adaptive algorithms in \citet[Sections 3.4--3.5]{doucet2012tutorial} make that reset explicit. If resampling is skipped, the previous unequal weights are retained. These conventions concern numerical particles, not changes to the company's contract.

For example, suppose two persistent contract types have current weights $3/4$ and $1/4$, and a new disclosure has likelihoods $1/3$ and $1$. Bayes' rule gives unnormalized masses $1/4$ and $1/4$, hence equal posterior probabilities. A four-child resampling allocation with three copies of the first type and one of the second gives aggregate incremental weights $3(1/3)=1$ and $1$, also equal. Multiplying by the old parent weights again would instead give $3/4$ and $1/4$. This exact allocation exposes the double counting; a random finite resampling draw need not reproduce the posterior exactly.

\citet[Sections 2--3]{doucet2012tutorial} explains sequential importance sampling, resampling and degeneracy. Its filtering setup conditions on a specified model. Static unknown parameters are not automatically handled well by adding them to the state and repeatedly resampling. That extension can suffer severe impoverishment.

For AccDSL, particle filtering is appropriate only after prior, transition and observation laws are admitted. It can be an optional numerical view of a history family. The family and source constraints remain available to detect support loss and distinguish probabilistic approximation from exact admissibility.
